\documentclass[11pt]{article}
\usepackage[margin=1in]{geometry}
\usepackage{amsmath,amssymb}
\usepackage{booktabs}
\usepackage{hyperref}
\usepackage{enumitem}
\hypersetup{colorlinks=true,linkcolor=blue,urlcolor=blue}

\title{Financial Turbulence Platform:\\
A Risk-Managed, Alpha-Tilted Multi-Asset Strategy}
\author{Shlok Parikh}
\date{June 2026}

\begin{document}
\maketitle

\begin{abstract}
We build a multi-asset allocation system that uses a statistical turbulence signal to scale
portfolio risk, a minimum-variance/vol-targeting base for stability, a leak-free exceedance
classifier as a leading de-risk trigger, and a time-series momentum tilt for return. Every
component is validated under a leakage-safe walk-forward backtest, net of transaction costs, on
ETF data spanning 2010--2026. A naively turbulence-tilted Black--Litterman strategy that is
\emph{dominated} by equal-weight (Sharpe $0.27$) is rebuilt into a strategy with Sharpe $0.82$,
Sortino $1.05$, and a maximum drawdown of $-8.8\%$ on an expanded 14-asset universe. The work
emphasises honest evaluation: a deep LSTM-CNN forecaster is shown to be statistically
indistinguishable from predicting the mean and is discarded in favour of a logistic regression
that genuinely beats persistence.
\end{abstract}

\section{Motivation}
Diversification fails precisely when it is most needed: asset correlations spike during crises, so
a portfolio that looks diversified in calm markets becomes concentrated in a single ``risk-on/off''
factor during a drawdown. The \emph{turbulence index} of Kritzman and Li (2010) measures this
directly, and the goal of this project is to turn it into a disciplined risk-scaling and allocation
system---and to do so with the methodological rigour (leakage control, honest baselines, robustness
testing) that separates a credible quant strategy from an over-fit backtest.

\section{Data and Universe}
The core universe is ten liquid ETFs spanning equities (SPY, EFA, EEM), real estate (VNQ), rates and
credit (AGG, LQD, HYG, TLT), and commodities (GLD, DBC), 2010--2026, daily. An \emph{expanded}
14-asset universe adds four distinct, low-correlation sleeves---T-bills (BIL), TIPS (TIP), oil (USO),
and EM bonds (EMB).

\paragraph{Effective breadth.} Eigenanalysis of the return correlation matrix shows the first
principal component (the common risk factor) explains $43.9\%$ of variance and the first three
explain $82\%$. The \emph{effective number of independent bets} (entropy of normalised eigenvalues)
is only $4.89$ out of $10$, and it \emph{collapses to $3.80$ during the COVID crash}---a quantitative
statement of why crisis diversification evaporates. Adding the four distinct sleeves raises effective
breadth to $6.54$, confirming they are genuinely diversifying rather than redundant.

\section{Methodology}
\subsection{Turbulence index}
For a return vector $x_t$ with rolling mean $\mu$ and covariance $\Sigma$ (estimated on a trailing
window, strictly in-sample), turbulence is the squared Mahalanobis distance
\begin{equation}
d_t^2 = (x_t-\mu)^\top \Sigma^{-1} (x_t-\mu),
\end{equation}
equivalently a one-sample Hotelling $T^2$ statistic. High $d_t^2$ means today's cross-asset move is
statistically improbable under the recent covariance regime.

\subsection{Allocation base}
Expected returns are the noisiest input to mean--variance optimisation, and inverting a raw sample
covariance amplifies estimation error (``error maximisation''). We therefore (i) shrink the
covariance with the Ledoit--Wolf estimator and (ii) use a \emph{risk-based} base---global
minimum-variance with a per-asset cap (risk parity and inverse-vol are also implemented)---which
needs no return forecast.

\subsection{Volatility targeting and the turbulence overlay}
The base book is scaled to a fixed annualised volatility target, then a turbulence \emph{exposure
overlay} scales gross exposure toward a cash sleeve (earning the risk-free rate) when the trailing
percentile of turbulence is elevated. This is the Kritzman--Li signal used as intended---to scale
\emph{risk}, not to rotate cross-sectionally into bonds.

\subsection{Forecasting: exceedance, not level}
Forecasting the turbulence \emph{level} is futile: its autocorrelation decays from $0.48$ (1d) to
$0.16$ (30d) to $-0.02$ (90d), so a level regressor merely re-derives persistence. A deep LSTM-CNN,
evaluated leak-free, does not beat a persistence baseline (7-day level rank-IC $0.146$ vs
persistence $0.175$) and is statistically indistinguishable from predicting the training mean. We
reframe the task as \emph{exceedance classification}---will turbulence cross its trailing 90th
percentile within $k$ days?---and find a regularised logistic regression beats persistence on
out-of-sample AUC (Table~\ref{tab:auc}); a gradient-boosted model over-fits. The deep model is
discarded.

\begin{table}[h]\centering
\caption{Out-of-sample exceedance AUC vs a persistence baseline.}\label{tab:auc}
\begin{tabular}{lccc}
\toprule
Horizon & Persistence & Logistic regression & Gradient boosting \\
\midrule
$k=5$  & 0.701 & \textbf{0.759} & 0.60 \\
$k=10$ & 0.685 & \textbf{0.727} & 0.50 \\
$k=21$ & 0.668 & \textbf{0.693} & 0.59 \\
\bottomrule
\end{tabular}
\end{table}

\subsection{Time-series momentum tilt}
Finally we add a return signal: time-series momentum (Moskowitz, Ooi \& Pedersen 2012). Each asset's
risk-adjusted 12-minus-1-month trend tilts the risk-based base (continuous tilt or a long-flat sign
filter). Beyond adding alpha, momentum is a \emph{principled} re-entry mechanism---it re-risks into
recoveries and shrinks persistent downtrends.

\subsection{Backtest protocol}
All results use a walk-forward backtest with monthly rebalancing, $10$ bps one-way costs, and a
strict no-lookahead guard (weights at rebalance $i$ use only data through $i{-}1$); the guard is unit
tested. Performance is summarised by annualised Sharpe and Sortino, annualised volatility, maximum
drawdown, and one-way turnover.

\section{Results}
\subsection{Building up the strategy}
Table~\ref{tab:ladder} isolates each layer's contribution on the core universe. The turbulence
overlay adds $+0.05$ Sharpe over naive vol-targeting; the forecast-driven overlay lifts Sortino from
$0.82$ to $0.92$ at lower turnover; momentum and the expanded universe compound to the best result.

\begin{table}[h]\centering
\caption{Incremental results (net of costs, leakage-safe). Top: core 10-asset universe,
full period. Bottom: expanded 14-asset universe.}\label{tab:ladder}
\begin{tabular}{lccccc}
\toprule
Configuration & Sharpe & Sortino & Ann.\ vol & Max DD & Turnover \\
\midrule
\multicolumn{6}{l}{\emph{Core universe (10 assets)}}\\
Equal-weight                 & 0.563 & 0.692 & 0.092 & $-0.218$ & 0.03 \\
Min-var + vol-target         & 0.644 & 0.767 & 0.051 & $-0.170$ & 0.43 \\
\quad + turbulence overlay   & 0.696 & 0.904 & 0.042 & $-0.160$ & 1.82 \\
\quad + momentum (sign)      & 0.740 & 0.943 & 0.042 & $-0.102$ & 2.62 \\
\midrule
\multicolumn{6}{l}{\emph{Expanded universe (14 assets)}}\\
Equal-weight                 & 0.443 & 0.532 & 0.086 & $-0.234$ & 0.03 \\
Strategy (no momentum)       & 0.714 & 0.912 & 0.025 & $-0.103$ & 1.83 \\
\textbf{+ momentum (continuous)} & \textbf{0.822} & \textbf{1.048} & \textbf{0.024} & $\mathbf{-0.088}$ & 2.02 \\
\bottomrule
\end{tabular}
\end{table}

\subsection{Marginal value of the forecast}
Over the common range 2013--2026, swapping the contemporaneous turbulence signal for the trailing
percentile of the forecast exceedance probability improves the overlay: Sharpe $0.63\!\to\!0.69$,
Sortino $0.82\!\to\!0.92$, max drawdown $-16\%\!\to\!-15\%$, with \emph{lower} turnover
($1.28$ vs $1.82$)---the gain is better timing, not more trading.

\subsection{Robustness}
The momentum tilt-strength sweep is smooth and monotonic (Sharpe $0.71\!\to\!0.83$ as strength goes
$0\!\to\!1$ on the expanded universe), the signature of a real effect rather than a fitted spike. By
sub-period (Sharpe), the continuous tilt is the more balanced choice: the sign filter helps the 2022
bear strongly but is weaker in 2016--2020, whereas the continuous tilt improves every sub-period
without large regressions.

\section{Crisis Stress Tests}
Table~\ref{tab:stress} reports the best config (expanded universe + continuous momentum) through six
stress episodes. The strategy cushions every crisis---most strikingly COVID ($-3.7\%$ drawdown vs SPY
$-35.7\%$) and the 2026 Iran-war oil shock ($-0.7\%$ vs SPY $-7.8\%$), a supply-driven dislocation it
was never tuned for, detected because turbulence keys off cross-asset correlation breakdown.

\begin{table}[h]\centering
\caption{Maximum drawdown within crisis windows (expanded universe + momentum).}\label{tab:stress}
\begin{tabular}{lcccc}
\toprule
Episode & SPY & Equal-weight & Strategy (base) & Strategy (+mom) \\
\midrule
2011 Euro crisis        & $-19.5\%$ & $-8.8\%$  & $-1.5\%$  & $-1.8\%$ \\
2015--16 China/oil      & $-13.2\%$ & $-10.9\%$ & $-2.0\%$  & $-2.0\%$ \\
2018 Q4 selloff         & $-19.8\%$ & $-9.1\%$  & $-1.1\%$  & $-1.0\%$ \\
COVID crash             & $-35.7\%$ & $-23.4\%$ & $-4.5\%$  & $-3.7\%$ \\
2022 bear (rates)       & $-26.2\%$ & $-16.6\%$ & $-9.4\%$  & $-8.3\%$ \\
2026 Iran war oil shock & $-7.8\%$  & $-3.0\%$  & $-0.9\%$  & $-0.7\%$ \\
\bottomrule
\end{tabular}
\end{table}

\section{Limitations}
We report these candidly, as they are part of the analysis:
\begin{itemize}[nosep]
  \item \textbf{Low absolute return.} The book is deliberately low-volatility ($\sim$2.4\%), so
  annualised return is modest ($\sim$2\%); the value is risk-adjusted, not headline return.
  \item \textbf{2022 is the hard regime.} When bonds and equities fall together, the
  minimum-variance base---which leans on bonds---offers less protection ($-8.3\%$).
  \item \textbf{Fast crashes cannot be fully dodged.} A 5--10 day signal cushions but does not avoid
  a two-week crash; the strategy also gives up upside in sharp recoveries.
  \item \textbf{Moderate statistical confidence.} One market history with roughly four stress
  episodes; cross-sectional breadth was added, but the number of independent time-series regimes is
  the binding constraint.
\end{itemize}

\section{Conclusion}
The system turns a broken, benchmark-dominated strategy (Sharpe $0.27$) into a rigorously validated,
risk-managed, alpha-tilted strategy (Sharpe $0.82$, Sortino $1.05$, max drawdown $-8.8\%$) that
demonstrably collapses crisis drawdowns. As much as the numbers, the contribution is the process:
leakage control, honest baselines (including discarding an over-fit deep model), robustness and
sub-period testing, and explicit limitations.

\begin{thebibliography}{9}
\bibitem{kl} M.~Kritzman and Y.~Li. Skulls, Financial Turbulence, and Risk Management.
\emph{Financial Analysts Journal}, 2010.
\bibitem{mop} T.~Moskowitz, Y.~H.~Ooi, and L.~H.~Pedersen. Time Series Momentum.
\emph{Journal of Financial Economics}, 2012.
\bibitem{lw} O.~Ledoit and M.~Wolf. Honey, I Shrunk the Sample Covariance Matrix.
\emph{Journal of Portfolio Management}, 2004.
\bibitem{bl} F.~Black and R.~Litterman. Global Portfolio Optimization.
\emph{Financial Analysts Journal}, 1992.
\end{thebibliography}

\end{document}
